%------------------------------------------------------------------------------%
\documentclass[fleqn,utf8,aspectratio=169,12pt]{beamer}

\usepackage{gaal}

\title{Volume do Paralelepípedo}

%------------------------------------------------------------------------------%
\begin{document}

\addtitle

%------------------------------------------------------------------------------%
\begin{frame}{Três Vetores Determinam um Paralelepípedo}
  Três vetores
  \(
    \vec{v}
  \),
  \(
    \vec{w}
  \) e
  \(
    \vec{z}
  \)
  com mesma origem

  \pause
  \bigskip
  Determinam um paralelepípedo no espaço $\R^3$

  \pause
  \bigskip
  Qual o volume desse paralelepípedo?
\end{frame}

%------------------------------------------------------------------------------%
\framedgraphic*{Três Vetores Determinam um Paralelepípedo}{E-11-Volume_paralelepipedo-1}{}
\framedgraphic*{Três Vetores Determinam um Paralelepípedo}{E-11-Volume_paralelepipedo-2}{}
\framedgraphic*{Três Vetores Determinam um Paralelepípedo}{E-11-Volume_paralelepipedo-3}{}

%------------------------------------------------------------------------------%
\begin{frame}{Volume}
  Volume do paralelepípedo
  \[
    \pause
    \text{Volume} = \text{Área da Base} \times \text{Altura}
  \]
  \[
    \pause
    V = Ah
  \]

  \pause
  A área da base é
  \[
    \pause
    A = \norm{\vec{v}\times\vec{w}}
  \]
\end{frame}

%------------------------------------------------------------------------------%
\framedgraphic*{Altura}{E-11-Volume_paralelepipedo-angulo-1}{}
\framedgraphic*{Altura}{E-11-Volume_paralelepipedo-angulo-2}{}
\framedgraphic*{Altura}{E-11-Volume_paralelepipedo-angulo-3}{}
\framedgraphic*{Altura}{E-11-Volume_paralelepipedo-angulo-4}{}
\framedgraphic*{Altura}{E-11-Volume_paralelepipedo-angulo-5}{}

%------------------------------------------------------------------------------%
\begin{frame}{Altura}
  Se $\phi$ é o ângulo entre $\vec{z}$ e o vetor normal à base, a altura é
  \[
    h = \norm{\vec{z}} \; \abs{\cos(\phi)}
  \]
\end{frame}

%------------------------------------------------------------------------------%
\begin{frame}{Volume}
  Volume
  \[
    V = \norm{\vec{v}\times\vec{w}} \; \norm{\vec{z}} \; \abs{\cos(\phi)}
  \]
  \pause
  Comparando com o produto escalar
  \[
    \vec{p} \cdot \vec{q} = \norm{\vec{p}} \; \norm{\vec{q}} \cos(\theta)
  \]
  \[
    \pause
    (\vec{v}\times\vec{w}) \cdot \vec{z} = \norm{\vec{v}\times\vec{w}} \; \norm{\vec{z}} \cos(\phi)
  \]
  \pause
  \[
    V = \abs{(\vec{v}\times\vec{w}) \cdot \vec{z}\,}
  \]
\end{frame}

%------------------------------------------------------------------------------%
\begin{frame}{Volume do Paralelepípedo}
  Dados três vetores
  \(
    \vec{v}
  \),
  \(
    \vec{w}
  \) e
  \(
    \vec{z}
  \)
  com mesma origem

  \pause
  \bigskip
  O volume do paralelepípedo é
  \[
    V \;=\;
    \abs{\,\left( \vec{v}\times\vec{w} \right)\cdot\vec{z}\,}
  \]

  \pause
  \smallskip
  O valor absoluto é necessário porque volume é uma grandeza não negativa
\end{frame}

%------------------------------------------------------------------------------%
%------------------------------------------------------------------------------%
\begin{question}
  Determine o volume do paralelepípedo determinado pelos vetores
  \[
    \vec{v} = \vetortri{1}{ 2}{0}
    \qquad
    \vec{w} = \vetortri{2}{-1}{1}
    \qquad
    \vec{z} = \vetortri{1}{ 0}{3}
  \]
\end{question}

%------------------------------------------------------------------------------%
\begin{resolution}{Resolução}
\begin{columns}[t]
\column{0.5\textwidth}
  \begin{align*}
    \onslide<+->{p & = \left( \vec{v}\times\vec{w} \right)\cdot\vec{z}}
    \\
    \onslide<+->{& = \begin{vmatrix}
      1 & 2  & 0 \\
      2 & -1 & 1 \\
      1 & 0  & 3
    \end{vmatrix}}
    \\
    \onslide<+->{& = \begin{vmatrix}
        -1 & 1 \\
        0  & 3
      \end{vmatrix}
      - 2 \begin{vmatrix}
        2 & 1 \\
        1 & 3
      \end{vmatrix}}
    \\
    \onslide<+->{& =- 3-10 \\}
    \onslide<+->{& = -13}
  \end{align*}
\column{0.5\textwidth}
  \[
    \onslide<+->{V}
    \onslide<+->{= \abs{p}}
    \onslide<+->{= \abs{-13}}
    \onslide<+->{= 13}
  \]
\end{columns}
\end{resolution}

%------------------------------------------------------------------------------%

%------------------------------------------------------------------------------%
\begin{question}
  Calcule o volume do paralelepípedo determinado pelos vetores
  \[
    \vec{v} = \vetortri{2}{ 1}{ 1}
    \qquad
    \vec{w} = \vetortri{3}{ 2}{-1}
    \qquad
    \vec{z} = \vetortri{1}{-1}{ 2}
  \]
\end{question}

%------------------------------------------------------------------------------%
\begin{resolution}{Resolução}
  \begin{align*}
    \onslide<+->{p & = \left( \vec{v}\times\vec{w} \right)\cdot\vec{z}}
    \onslide<+->{\;=\; \begin{vmatrix}
                         2 & 1  & 1  \\
                         3 & 2  & -1 \\
                         1 & -1 & 2
                       \end{vmatrix}}
    \\
    \onslide<+->{& = 2
      \begin{vmatrix}
        2  & -1 \\
        -1 & 2
      \end{vmatrix}
      -
      \begin{vmatrix}
        3 & -1 \\
        1 & 2
      \end{vmatrix}
      +
      \begin{vmatrix}
        3 & 2  \\
        1 & -1
      \end{vmatrix}}
    \\
    \onslide<+->{& = 2(4-1) - (6+1) + (-3-2)} \\
    \onslide<+->{& = 2 \times 3 - 7 - 5}
    \onslide<+->{\;=\; -6}
  \end{align*}

  \[
    \onslide<+->{V}
    \onslide<+->{\;=\; \abs{p}}
    \onslide<+->{\;=\; \abs{-6}}
    \onslide<+->{\;=\; 6}
  \]
\end{resolution}

%------------------------------------------------------------------------------%

%------------------------------------------------------------------------------%
\begin{question}
  Considere os pontos
  \begin{align*}
    A & = ( 1, 1, 0 ) \\
    B & = ( 3, 2, 1 ) \\
    C & = ( 2, 4, 1 ) \\
    D & = ( 1, 1, 3 )
  \end{align*}
  Determine o volume do paralelepípedo que possui $A$ como vértice
  e é determinado pelos vetores
  \[
    \overrightarrow{AB}
    \qquad
    \overrightarrow{AC}
    \qquad
    \overrightarrow{AD}
  \]
\end{question}

%------------------------------------------------------------------------------%
\begin{resolution}{Resolução}
  \[
    \overrightarrow{AB}
    \pause
    = B-A
    \pause
    = \vetortri{2}{1}{1}
  \]

  \[
    \pause
    \overrightarrow{AC}
    \pause
    = C-A
    \pause
    = \vetortri{1}{3}{1}
  \]

  \[
    \pause
    \overrightarrow{AD}
    \pause
    = D-A
    \pause
    = \vetortri{0}{0}{3}
  \]
\end{resolution}

%------------------------------------------------------------------------------%
\begin{resolution}{Resolução}
\begin{columns}[t]
\column{0.5\textwidth}
  \begin{align*}
    \onslide<+->{p}
    \onslide<+->{& = \left(\overrightarrow{AB} \times \overrightarrow{AC}\right) \cdot \overrightarrow{AD}}
    \\
    \onslide<+->{& = \begin{vmatrix}
          2 & 1 & 1 \\
          1 & 3 & 1 \\
          0 & 0 & 3
        \end{vmatrix}}
    \\
    \onslide<+->{& = 3 (-1)^{3+3}
        \begin{vmatrix}
          2 & 1 \\
          1 & 3
        \end{vmatrix}}
    \\
    \onslide<+->{& = 3(6-1)} \\
    \onslide<+->{& = 15}
  \end{align*}
\column{0.5\textwidth}
  \[
    \onslide<+->{V}
    \onslide<+->{= \abs{p}}
    \onslide<+->{= \abs{15}}
    \onslide<+->{= 15}
  \]
\end{columns}
\end{resolution}
%------------------------------------------------------------------------------%

%%------------------------------------------------------------------------------%
\begin{question}
  Considere os pontos
  \[
  A=(0,0,0)
  \]
  \[
  B=(2,1,0)
  \]
  \[
  C=(1,3,2)
  \]
  \[
  D=(0,1,4)
  \]
  Determine o volume do tetraedro $ABCD$
\end{question}

%------------------------------------------------------------------------------%
\begin{resolution}{Exemplo 4 — Resolução}
Construímos
\[
\overrightarrow{AB}=(2,1,0)
\]
\[
\overrightarrow{AC}=(1,3,2)
\]
\[
\overrightarrow{AD}=(0,1,4)
\]
O paralelepípedo correspondente possui volume
\[
V_P=
\left|
\begin{vmatrix}
2&1&0\\
1&3&2\\
0&1&4
\end{vmatrix}
\right|
\]
\end{resolution}

%------------------------------------------------------------------------------%
\begin{resolution}{Exemplo 4 — Resolução}
Calculando o determinante
\[
V_P=
\left|
2
\begin{vmatrix}
3&2\\
1&4
\end{vmatrix}
-
\begin{vmatrix}
1&2\\
0&4
\end{vmatrix}
\right|
\]
\[
V_P=|2(10)-4|
\]
\[
V_P=16
\]
Um tetraedro ocupa um sexto do paralelepípedo determinado pelos mesmos três vetores
\[
V_T=\frac16V_P
\]
\[
V_T=\frac{16}{6}
\]
\[
V_T=\frac83
\]
\end{resolution}
%------------------------------------------------------------------------------%

%%------------------------------------------------------------------------------%
\begin{question}
  Determine $k$ para que o paralelepípedo determinado por
  \[
  \vec{v}=(1,2,1)
  \]
  \[
  \vec{w}=(2,-1,3)
  \]
  \[
  \vec{z}=(k,1,2)
  \]
  tenha volume igual a $6$
\end{question}

%------------------------------------------------------------------------------%
\begin{resolution}{Exemplo 5 — Resolução}
O volume é
\[
V=
\left|
\begin{vmatrix}
1&2&1\\
2&-1&3\\
k&1&2
\end{vmatrix}
\right|
\]
Expandindo pela primeira linha
\[
V=
\left|
\begin{vmatrix}
-1&3\\
1&2
\end{vmatrix}
-
2
\begin{vmatrix}
2&3\\
k&2
\end{vmatrix}
+
\begin{vmatrix}
2&-1\\
k&1
\end{vmatrix}
\right|
\]
\[
V=
|-5-2(4-3k)+(2+k)|
\]
\[
V=
|6k-11|
\]
Como o volume deve ser $6$
\[
|6k-11|=6
\]
\end{resolution}

%------------------------------------------------------------------------------%
\begin{resolution}{Exemplo 5 — Resolução}
Temos duas possibilidades
\[
6k-11=6
\]
ou
\[
6k-11=-6
\]
Da primeira
\[
6k=17
\]
\[
k=\frac{17}{6}
\]
Da segunda
\[
6k=5
\]
\[
k=\frac56
\]
Portanto os valores possíveis são
\[
k=\frac56
\]
e
\[
k=\frac{17}{6}
\]
\end{resolution}
%------------------------------------------------------------------------------%


%%------------------------------------------------------------------------------%
%\section{Lista Mínima}
%
%%------------------------------------------------------------------------------%
%\begin{frame}{Lista Mínima}
%  \begin{enumerate}
%    \item
%  \end{enumerate}
%
%  \vfill
%  Atenção: A prova é baseada no livro, não nas apresentações
%\end{frame}

%------------------------------------------------------------------------------%
\end{document}
%------------------------------------------------------------------------------%
